%%
%% Test: Box Styles — v1.1
%% Stage: boxes refactor (stages 1-4)
%% 
%% Purpose: Comprehensive test of all box types:
%%   - All 10 environments (theorem, lemma, corollary, proposition,
%%     definition, example, remark, exercise, solution, proof)
%%   - Short and long boxes
%%   - Breakable boxes (span multiple pages)
%%   - Boxes with formulas, lists, code
%%   - Boxes on odd/even pages
%% 

\documentclass{matinbook}

\title{تست جامع جعبه‌ها — نسخه ۱.۱}
\author{تیم MatinBook}
\date{۱۴۰۵}

\booktitle{تست جامع جعبه‌ها}

\begin{document}

\maketitle

\tableofcontents

\chapter{تست جامع جعبه‌های رنگی}
\label{chap:test-boxes}

این سند، تمام جعبه‌های رنگی \texttt{mb-boxes} را به‌صورت
جامع بررسی می‌کند. هدف این است که ببینیم آیا همه‌ی محیط‌ها
با استایل‌های مختلف به‌درستی کار می‌کنند.

%========================================
\section{خانواده‌ی قضیه (آبی)}
\label{sec:theorem-family}
%========================================

\subsection{قضیه — کوتاه}
\begin{theorem}[قضیه‌ی فیثاغورس]
    در هر مثلث قائم‌الزاویه، مربع وتر برابر است با مجموع
    مربعات دو ضلع دیگر:
    \[
    a^2 + b^2 = c^2
    \]
    \label{thm:pythagoras}
\end{theorem}

\subsection{لم — کوتاه}
\begin{lemma}[لم اقلیدس]
    اگر $p$ عددی اول باشد و $p \mid ab$، آنگاه $p \mid a$
    یا $p \mid b$.
    \label{lem:euclid}
\end{lemma}

\subsection{نتیجه — کوتاه}
\begin{corollary}
    اگر $p$ عددی اول باشد و $p \mid a_1 a_2 \cdots a_n$،
    آنگاه $p$ حداقل یکی از $a_i$ها را می‌شمارد.
    \label{cor:euclid-general}
\end{corollary}

\subsection{گزاره — کوتاه}
\begin{proposition}
    کوچک‌ترین مقسوم‌علیه بزرگ‌تر از ۱ هر عدد طبیعی $n > 1$،
    یک عدد اول است.
    \label{prop:smallest-divisor}
\end{proposition}

%========================================
\section{تعریف}
\label{sec:definition}
%========================================

\begin{definition}[عدد اول]
    عدد طبیعی بزرگ‌تر از ۱ که تنها دو مقسوم‌علیه مثبت دارد،
    \emph{عدد اول} نامیده می‌شود. اعدادی که اول نباشند را
    \emph{مرکب} گوییم.
    \label{def:prime}
\end{definition}

\begin{definition}[گروه]
    یک \textbf{گروه} $(G, *)$ شامل یک مجموعه‌ی ناتهی $G$
    و یک عمل دوتایی $*$ روی $G$ است که در چهار شرط زیر
    صدق کند:
    \begin{enumerate}
        \item \textbf{بسته بودن:} $\forall a, b \in G: a * b \in G$
        \item \textbf{شرکت‌پذیری:} $\forall a, b, c \in G: (a * b) * c = a * (b * c)$
        \item \textbf{عضو خنثی:} $\exists e \in G: \forall a \in G: e * a = a * e = a$
        \item \textbf{عضو وارون:} $\forall a \in G: \exists a^{-1} \in G: a * a^{-1} = a^{-1} * a = e$
    \end{enumerate}
    \label{def:group}
\end{definition}

%========================================
\section{مثال و نکته}
\label{sec:example-remark}
%========================================

\begin{example}[اعداد اول کوچک]
    اعداد اول کوچک‌تر از ۳۰ عبارتند از:
    \[
    2, 3, 5, 7, 11, 13, 17, 19, 23, 29
    \]
    \label{ex:small-primes}
\end{example}

\begin{remark}[نکته]
    عدد ۱ نه اول است و نه مرکب. این قرارداد برای سادگی
    در بیان قضایا اتخاذ شده است.
    \label{rem:one-not-prime}
\end{remark}

\begin{remark}
    عدد ۲ تنها عدد اول زوج است. تمام اعداد اول دیگر فرد
    هستند.
    \label{rem:two-is-special}
\end{remark}

%========================================
\section{تمرین}
\label{sec:exercise}
%========================================

\begin{exercise}[مجموع اعداد زوج]
    ثابت کنید که مجموع دو عدد زوج، زوج است.
    \label{exe:sum-even}
\end{exercise}

\begin{exercise}
    ثابت کنید که $(\Z_n, +_n)$ یک گروه است.
    \label{exe:zn-group}
\end{exercise}

%========================================
\section{راه‌حل و اثبات}
\label{sec:solution-proof}
%========================================

\begin{proof}
    فرض کنید $a = 2m$ و $b = 2n$. آنگاه:
    \[
    a + b = 2m + 2n = 2(m + n)
    \]
    که زوج است.
\end{proof}

\begin{solution}
    برای اثبات این که $(\Z_n, +_n)$ یک گروه است، چهار شرط
    گروه را بررسی می‌کنیم:
    \begin{enumerate}
        \item \textbf{بسته بودن:} مجموع دو عدد به پیمانه $n$
        همیشه در $\Z_n$ است.
        \item \textbf{شرکت‌پذیری:} جمع به پیمانه $n$ شرکت‌پذیر است.
        \item \textbf{عضو خنثی:} عدد $0$ است.
        \item \textbf{عضو وارون:} وارون $a$، عدد $n - a$ است.
    \end{enumerate}
\end{solution}

%========================================
\section{قضیه‌ی طولانی (تست شکست)}
\label{sec:long-theorem}
%========================================

این بخش برای بررسی شکست باکس \texttt{matinbox} در صفحات
طراحی شده است.

\begin{theorem}[قضیه‌ی اساسی حساب]
    هر عدد طبیعی $n > 1$ را می‌توان به‌طور یکتا (صرف‌نظر از
    ترتیب عوامل) به‌صورت حاصل‌ضرب اعداد اول نوشت:
    \[
    n = p_1^{e_1} \cdot p_2^{e_2} \cdots p_k^{e_k}
    \]
    که در آن $p_1 < p_2 < \cdots < p_k$ اعداد اول و
    $e_1, e_2, \ldots, e_k \geq 1$ اعداد صحیح مثبت هستند.

    این قضیه یکی از مهم‌ترین قضایای نظریه‌ی اعداد است.
    در ادامه، چند نکته درباره‌ی آن بیان می‌شود.

    \textbf{نکته‌ی اول — یکتایی:} منظور از یکتایی این است
    که صرف‌نظر از ترتیب عوامل، تجزیه همیشه یکسان است.
    برای مثال، عدد ۶۰ را می‌توان به‌صورت زیر نوشت:
    \[
    60 = 2 \cdot 2 \cdot 3 \cdot 5 = 2^2 \cdot 3 \cdot 5
    \]

    \textbf{نکته‌ی دوم — کاربرد در رمزنگاری:} قضیه‌ی اساسی
    حساب پایه‌ی الگوریتم RSA است. در این الگوریتم، دو عدد
    اول بزرگ $p$ و $q$ انتخاب می‌شوند و $n = pq$ محاسبه
    می‌شود. امنیت RSA بر این فرض استوار است که تجزیه‌ی
    $n$ به عوامل اول برای $n$های بزرگ بسیار دشوار است.

    \textbf{نکته‌ی سوم — تعمیم:} قضیه‌ی اساسی حساب در
    حلقه‌ی اعداد صحیح $\Z$ برقرار است، ولی در حلقه‌های
    دیگر ممکن است برقرار نباشد. برای مثال، در حلقه‌ی
    $\Z[\sqrt{-5}]$، عدد ۶ دو تجزیه‌ی متفاوت دارد:
    \[
    6 = 2 \cdot 3 = (1 + \sqrt{-5})(1 - \sqrt{-5})
    \]

    \textbf{نکته‌ی چهارم — تاریخچه:} این قضیه اولین بار
    توسط اقلیدس در کتاب هفتم «اصول» به‌صورت ضمنی بیان شد.

    \textbf{نکته‌ی پنجم — مثال‌های بیشتر:}
    \begin{itemize}
        \item $12 = 2^2 \cdot 3$
        \item $30 = 2 \cdot 3 \cdot 5$
        \item $100 = 2^2 \cdot 5^2$
        \item $210 = 2 \cdot 3 \cdot 5 \cdot 7$
        \item $1024 = 2^{10}$
    \end{itemize}

    \textbf{نکته‌ی ششم — الگوریتم تجزیه:} برای یافتن
    تجزیه‌ی یک عدد، می‌توان از الگوریتم تقسیم مکرر
    استفاده کرد.

    \textbf{نکته‌ی هفتم — پیچیدگی محاسباتی:} مسئله‌ی
    تجزیه‌ی اعداد بزرگ به عوامل اول، یکی از مسائل مهم
    در نظریه‌ی پیچیدگی محاسباتی است.

    \textbf{نکته‌ی هشتم — کاربردهای دیگر:} علاوه بر
    رمزنگاری، قضیه‌ی اساسی حساب در بسیاری از شاخه‌های
    دیگر ریاضیات کاربرد دارد.

    \textbf{نکته‌ی نهم — نتیجه‌گیری:} بنابراین، قضیه‌ی
    اساسی حساب یکی از پایه‌ای‌ترین قضایای ریاضیات است
    که در سراسر نظریه‌ی اعداد، جبر مجرد، و رمزنگاری
    کاربرد دارد.

    \label{thm:fundamental-arithmetic}
\end{theorem}

%========================================
\section{اثبات طولانی (تست شکست outline)}
\label{sec:long-proof}
%========================================

این بخش برای بررسی شکست باکس \texttt{matin-outline} طراحی
شده است.

\begin{proof}[اثبات بی‌نهایت بودن اعداد اول]
    می‌خواهیم ثابت کنیم که بی‌نهایت عدد اول وجود دارد.
    اثبات را با برهان خلف انجام می‌دهیم.

    فرض کنید تعداد اعداد اول متناهی باشد. آن‌ها را
    $p_1, p_2, \ldots, p_n$ بنامید. عدد زیر را در نظر بگیرید:
    \[
    N = p_1 \cdot p_2 \cdots p_n + 1
    \]

    \textbf{حالت اول:} اگر $N$ خودش اول باشد، آنگاه عدد
    اولی یافته‌ایم که در فهرست $p_1, \ldots, p_n$ نیست.
    این با فرض متناهی بودن اعداد اول تناقض دارد.

    \textbf{حالت دوم:} اگر $N$ مرکب باشد، آنگاه عامل اولی
    مانند $q$ دارد. این $q$ باید یکی از $p_1, \ldots, p_n$
    باشد. اما $N$ بر هیچ‌کدام از $p_i$ها بخش‌پذیر نیست:
    \[
    N \bmod p_i = (p_1 \cdots p_n + 1) \bmod p_i = 1
    \]

    در هر دو حالت به تناقض رسیدیم. پس فرض متناهی بودن
    اعداد اول باطل است.

    \textbf{نکته:} این اثبات اولین بار توسط اقلیدس ارائه شد
    و یکی از زیباترین اثبات‌های ریاضیات است.

    \textbf{اثبات جایگزین اویلر:}
    \[
    \sum_{p \text{ اول}} \frac{1}{p} = \infty
    \]

    \textbf{قضیه‌ی دیریکله:} در هر تصاعد حسابی به شکل
    $a, a+d, a+2d, \ldots$ که $\gcd(a, d) = 1$،
    بی‌نهایت عدد اول وجود دارد.

    \textbf{کاربرد در رمزنگاری:} نتیجه‌ی بی‌نهایت بودن
    اعداد اول، پایه‌ی بسیاری از الگوریتم‌های رمزنگاری
    مدرن است.

    \textbf{فاصله‌ی بین اعداد اول:} این اثبات نشان می‌دهد
    که فاصله‌ی بین اعداد اول می‌تواند به‌طور دلخواه بزرگ
    باشد.

    \textbf{حدس گلدباخ:} هر عدد زوج بزرگ‌تر از ۲ را می‌توان
    به‌صورت مجموع دو عدد اول نوشت.

    \textbf{حدس اعداد اول دوقلو:} بی‌نهایت جفت عدد اول
    به شکل $(p, p+2)$ وجود دارد.

    \textbf{نتیجه‌ی نهایی:} بنابراین، بی‌نهایت عدد اول
    وجود دارد.
        می‌خواهیم ثابت کنیم که بی‌نهایت عدد اول وجود دارد.
    اثبات را با برهان خلف انجام می‌دهیم.

    فرض کنید تعداد اعداد اول متناهی باشد. آن‌ها را
    $p_1, p_2, \ldots, p_n$ بنامید. عدد زیر را در نظر بگیرید:
    \[
    N = p_1 \cdot p_2 \cdots p_n + 1
    \]

    \textbf{حالت اول:} اگر $N$ خودش اول باشد، آنگاه عدد
    اولی یافته‌ایم که در فهرست $p_1, \ldots, p_n$ نیست.
    این با فرض متناهی بودن اعداد اول تناقض دارد.

    \textbf{حالت دوم:} اگر $N$ مرکب باشد، آنگاه عامل اولی
    مانند $q$ دارد. این $q$ باید یکی از $p_1, \ldots, p_n$
    باشد. اما $N$ بر هیچ‌کدام از $p_i$ها بخش‌پذیر نیست:
    \[
    N \bmod p_i = (p_1 \cdots p_n + 1) \bmod p_i = 1
    \]

    در هر دو حالت به تناقض رسیدیم. پس فرض متناهی بودن
    اعداد اول باطل است.

    \textbf{نکته:} این اثبات اولین بار توسط اقلیدس ارائه شد
    و یکی از زیباترین اثبات‌های ریاضیات است.

    \textbf{اثبات جایگزین اویلر:}
    \[
    \sum_{p \text{ اول}} \frac{1}{p} = \infty
    \]

    \textbf{قضیه‌ی دیریکله:} در هر تصاعد حسابی به شکل
    $a, a+d, a+2d, \ldots$ که $\gcd(a, d) = 1$،
    بی‌نهایت عدد اول وجود دارد.

    \textbf{کاربرد در رمزنگاری:} نتیجه‌ی بی‌نهایت بودن
    اعداد اول، پایه‌ی بسیاری از الگوریتم‌های رمزنگاری
    مدرن است.

    \textbf{فاصله‌ی بین اعداد اول:} این اثبات نشان می‌دهد
    که فاصله‌ی بین اعداد اول می‌تواند به‌طور دلخواه بزرگ
    باشد.

    \textbf{حدس گلدباخ:} هر عدد زوج بزرگ‌تر از ۲ را می‌توان
    به‌صورت مجموع دو عدد اول نوشت.

    \textbf{حدس اعداد اول دوقلو:} بی‌نهایت جفت عدد اول
    به شکل $(p, p+2)$ وجود دارد.

    \textbf{نتیجه‌ی نهایی:} بنابراین، بی‌نهایت عدد اول
    وجود دارد.
\end{proof}

%========================================
\section{باکس با لیست}
\label{sec:boxes-with-lists}
%========================================

\begin{example}[لیست داخل مثال]
    داخل این مثال، یک لیست داریم:
    \begin{itemize}
        \item مورد اول
        \item مورد دوم
        \item مورد سوم
    \end{itemize}
    و یک لیست شماره‌دار:
    \begin{enumerate}
        \item اول
        \item دوم
        \item سوم
    \end{enumerate}
\end{example}

\begin{theorem}[قضیه با چند بخش]
    \textbf{بخش اول:} این یک متن آزمایشی است.

    \textbf{بخش دوم:} این هم یک متن دیگر است.

    \textbf{بخش سوم:} و این هم بخش سوم.
\end{theorem}

%========================================
\section{باکس با فرمول‌های مختلف}
\label{sec:boxes-with-formulas}
%========================================

\begin{example}[فرمول‌های مختلف]
    فرمول درون‌خطی: $E = mc^2$.

    فرمول جدا:
    \[
    \int_{-\infty}^{+\infty} e^{-x^2} \, dx = \sqrt{\pi}
    \]

    فرمول چندخطی:
    \begin{align*}
        \sum_{n=1}^{\infty} \frac{1}{n^2} &= \frac{\pi^2}{6} \\
        \lim_{n \to \infty} \left(1 + \frac{1}{n}\right)^n &= e
    \end{align*}
\end{example}

%========================================
\section{ارجاعات}
\label{sec:references}
%========================================

ارجاع به قضیه‌ی فیثاغورس: \cref{thm:pythagoras}

ارجاع به لم اقلیدس: \cref{lem:euclid}

ارجاع به نتیجه: \cref{cor:euclid-general}

ارجاع به گزاره: \cref{prop:smallest-divisor}

ارجاع به تعریف عدد اول: \cref{def:prime}

ارجاع به تعریف گروه: \cref{def:group}

ارجاع به مثال: \cref{ex:small-primes}

ارجاع به نکته: \cref{rem:one-not-prime}

ارجاع به تمرین: \cref{exe:sum-even}

ارجاع به قضیه‌ی اساسی حساب: \cref{thm:fundamental-arithmetic}

%========================================
\section{نتیجه‌گیری}
\label{sec:conclusion}
%========================================

اگر این سند بدون خطا کامپایل شود، یعنی:

محیط \texttt{theorem} با استایل آبی و نوار عنوان کار می‌کند.

محیط \texttt{lemma} با استایل آبی تیره‌تر کار می‌کند.

محیط \texttt{corollary} با استایل آبی کار می‌کند.

محیط \texttt{proposition} با استایل آبی کار می‌کند.

محیط \texttt{definition} با استایل سبز کار می‌کند.

محیط \texttt{example} با استایل نارنجی outline کار می‌کند.

محیط \texttt{remark} با استایل قرمز کار می‌کند.

محیط \texttt{exercise} با استایل بنفش کار می‌کند.

محیط \texttt{solution} با استایل سبز outline کار می‌کند.

محیط \texttt{proof} با استایل خاکستری outline کار می‌کند.

باکس‌های طولانی در صفحات می‌شکنند.

باکس‌های با لیست و فرمول درست کار می‌کنند.

ارجاعات \texttt{cleveref} درست کار می‌کنند.

\textbf{وضعیت تست:} قبول.

\end{document}